\section{Введение}
\label{sec:Chapter0} \index{Chapter0}

\subsection{Актуальность}

Трассировка лучей (ray tracing) в последнее десятилетие стала
основным алгоритмом физически корректного рендеринга и активно
проникает в реальное время благодаря появлению специализированных
аппаратных блоков NVIDIA RT~Cores и AMD Ray~Accelerators, а также
программных интерфейсов DirectX~Raytracing (DXR) и Vulkan
\texttt{KHR\_ray\_tracing}. Однако стоимость одного ray-primitive
теста растёт линейно с числом примитивов сцены, поэтому без
ускоряющей пространственной структуры данный подход
непригоден для интерактивных применений~\cite{paper}.

Стандартом де-факто для таких структур стала иерархия охватывающих
объёмов (Bounding~Volume~Hierarchy, BVH). При трассировке луч
обходит дерево BVH и отбрасывает поддеревья, AABB которых не
пересекается с лучом, что переводит сложность поиска ближайшего
пересечения в среднем к $O(\log n)$. Качество BVH принято измерять
поверхностно-площадной эвристикой (Surface Area Heuristic, SAH)~\cite{paper, ganestam16_splitting}:
\begin{equation}
    c(N) = c_T \sum_{N_i \in \text{Internal}} \frac{SA(N_i)}{SA(N)} + c_I \sum_{N_l \in \text{Leaf}} \frac{SA(N_l)}{SA(N)}\,|N_l|
\end{equation}
где:
\begin{itemize}
    \item $N$ --- корневой узел поддерева;
    \item $\text{Internal}$ и $\text{Leaf}$ --- множества внутренних (развилки дерева, хранящие AABB потомков) и листовых узлов (хранящие сами треугольники);
    \item $SA(X)$ --- площадь поверхности AABB узла $X$;
    \item $c_T$ и $c_I$ --- константные стоимости обхода узла и теста пересечения с треугольником;
    \item $|N_l|$ --- число треугольников в листе $N_l$.
\end{itemize}

Формула считает ожидаемую стоимость трассировки луча на основе геометрической вероятности: отношение площадей $SA(X) / SA(N)$ --- это шанс того, что луч, прошедший через $N$, пересечёт и вложенный узел $X$. При этом, если AABB дочерних узлов сильно перекрываются, их суммарная площадь увеличивается, лучу чаще приходится обходить обе ветви, и рендеринг становится дороже.

В современных графических приложениях BVH строится в два уровня:
\begin{itemize}
    \item \textbf{BLAS (Bottom-Level Acceleration Structure)} --- структуры нижнего уровня. Каждая из них содержит вершины и индексы треугольников отдельного меша или 3D-модели в локальной системе координат.
    \item \textbf{TLAS (Top-Level Acceleration Structure)} --- структура верхнего уровня. Она не содержит ассетов напрямую, а состоит из инстансов BLAS, расставленных по сцене с определёнными матрицами трансформаций (позиция, вращение, масштаб).
\end{itemize}
Такое разделение позволяет быстро обновлять сцену при движении объектов (достаточно перестроить только TLAS), не перестраивая тяжёлые BLAS для самих ассетов. Общая схема двухуровневой структуры и место, в котором возникает исследуемая проблема, показаны на приведённой ниже схеме.

\begin{figure}[h!]
\centering
\includegraphics[width=0.82\textwidth]{diagrams/scheme_tlas_blas.png}
\caption{Двухуровневая структура ускорения в семпле \texttt{testGI}:
сцена из 16 инстансов модели порождает TLAS, листья которого ссылаются
на BLAS отдельных ассетов. Луч прохода \texttt{rtsm::do\_trace} сначала
обходит TLAS, затем спускается в BLAS. Именно качество BLAS и степень
перекрытия их AABB в TLAS определяют стоимость трассировки.}
\label{fig:arch_tlas_blas}
\end{figure}

На этом фоне построение BLAS, как правило, выполняется драйвером графического
вендора и недоступно для прямого вмешательства разработчиком: API
DXR/Vulkan~RT принимает на вход \emph{набор ассетов} (vertex
buffer + index buffer + диапазон), а внутреннюю стратегию построения
выбирает драйвер. Это означает, что основной канал влияния на
качество BLAS со стороны разработчика --- \textbf{предобработка
меша до его передачи билдеру}.

\subsection{Объект и предмет исследования}

\textbf{Объект исследования} --- процесс построения и обхода ускоряющих структур данных для аппаратной трассировки лучей в реальном времени.

\textbf{Предмет исследования} --- методы и алгоритмы пространственной кластеризации и предобработки ассетов трёхмерных моделей, направленные на повышение качества структур нижнего уровня и ускорение рендеринга.

\subsection{Ключевые понятия}

\defitem{BVH}{Bounding~Volume~Hierarchy. Двоичное дерево, в листьях
  которого хранятся примитивы сцены, а во внутренних узлах --
  AABB всех потомков.}
\defitem{AABB}{Axis-Aligned Bounding Box. Параллелепипед со сторонами,
  параллельными осям координат, содержащий поддерево примитивов.}
\defitem{BLAS / TLAS}{Bottom- и Top-Level Acceleration Structure --
  два уровня иерархии BVH в современных графических API (DXR / Vulkan~RT). BLAS хранит сами ассеты модели, а TLAS организует их экземпляры на сцене.}
\defitem{RElem}{Атомарная меш-структура исследуемого движка
  DagorEngine (\texttt{ShaderMesh::RElem}). Содержит ссылку на
  материал и непрерывный поддиапазон вершин/индексов общего
  буфера; одному \texttt{RElem} соответствует одна
  \texttt{RaytraceGeometryDescription} на входе
  билдера.}
\defitem{Morton-код}{Числовая характеристика 3D-точки, полученная
  чередованием битов её нормированных координат. Сортировка
  точек по Morton-коду эквивалентна обходу пространства по
  заполняющей кривой Z-order, что обеспечивает локальность
  соседних элементов в массиве.}
\defitem{SAH cost}{Поверхностно-площадная стоимость BVH;
  общепринятая офлайн-мера качества дерева,
  коррелирующая с реальным временем
  трассировки~\cite{paper}.}
\defitem{RTSM}{Ray~Traced~Shadow~Map. Подсистема трассировки теней
  в DagorEngine, выполняющая ray-cast от каждого пикселя экрана
  в сторону источника света; целевая нагрузка для измерения
  ускорения.}

\subsection{Постановочная гипотеза}

Графический конвейер часто выделяет отдельный меш-ассет и, соответственно, отдельный BLAS под каждый уникальный материал, например, дерево или металл. Если элементы такого материала разбросаны по всему объёму крупной модели, то AABB соответствующего BLAS будет сопоставим с размерами всей модели. В результате в TLAS возникает множество огромных и полностью перекрывающихся инстансов --- по одному на каждый материал. При сборке TLAS AABB его листовых узлов получаются огромными и сильно перекрывающимися друг с другом. При трассировке луча алгоритм обхода вынужден проверять пересечение со всеми этими крупными BLAS, что может существенно снижать производительность рендеринга.

Изложенное обосновывает следующую гипотезу работы: \emph{объединение всех треугольников модели с последующим разделением её на множество независимых, пространственно-компактных мини-ассетов, выполненное на стороне CPU до подачи данных в билдер, демонстрирует устойчивое снижение времени аппаратной GPU-трассировки лучей на сценах с таким типом ассетов.}

\subsection{Степень разработанности темы}

Проблема качества BVH активно исследуется: классические top-down SAH-билдеры,
HLBVH и SBVH повышают качество дерева при пространственно компактном входе,
однако в DXR/Vulkan~RT построение BLAS выполняется драйвером и недоступно
для модификации изнутри. Отдельные приёмы предобработки (Morton-reorder,
SAH pre-splitting) совместимы с API, но не устраняют перекрытия в TLAS,
возникающие при материальном разбиении крупных моделей. Промышленные
руководства NVIDIA и AMD рекомендуют вручную нарезать «жирные» ассеты,
не предлагая автоматизированного runtime-решения. Таким образом, автоматическая
пространственная кластеризация на входе вендорского билдера остаётся
малоосвоенной.

\subsection{Цель и задачи}

\textbf{Цель} --- разработать и экспериментально оценить алгоритм
рантайм-предобработки мешей, повышающий качество BLAS, построенного
вендорским билдером, и достичь измеримого ускорения трассировки
лучей в реальной сцене игрового движка без визуальных потерь.

\textbf{Задачи}:
\begin{enumerate}
    \item Провести обзор существующих подходов к улучшению качества
      BVH, оценить их применимость в
      условиях фиксированного DXR-pipeline.
    \item Изучить цепочку формирования BLAS в DagorEngine и
      определить точку, в которой возможна предобработка меша без
      модификации драйвера.
    \item Реализовать алгоритм глобальной пространственной кластеризации ассетов на
      стороне runtime: объединение всех треугольников модели в
      единый массив, сортировка по Morton-коду, разбиение на $K$
      пространственно-компактных кластеров с собственными
      vertex/index-буферами.
    \item Провести замер ускорения трассировки на сцене с
      «плохими» ассетами в семпле \texttt{testGI} и подтвердить
      визуальную корректность результата.
    \item Сформулировать ограничения метода и направления
      дальнейшей работы.
\end{enumerate}

\subsection{Методы исследования}

Для решения поставленных задач использованы следующие методы.
\textbf{Теоретический} --- обзор и сравнительный анализ подходов к построению
и улучшению BVH с позиции совместимости с DXR-pipeline; формализация входных данных, ограничений и
критериев оценки.
\textbf{Проектный} --- анализ цепочки формирования BLAS в DagorEngine,
выделение точки вмешательства и проектирование алгоритма глобальной
пространственной кластеризации.
\textbf{Экспериментальный} --- программная реализация метода в семпле
\texttt{testGI}, измерение медианного GPU-времени \texttt{rtsm::do\_trace}
и качественная проверка пространственной локализации BLAS в Nsight~Graphics. Выбор GPU-замера вместо офлайн-метрик SAH
обоснован тем, что в DXR реальная стоимость трассировки определяется
драйверо-специфичным билдером и аппаратным обходом.

\subsection{Научная новизна}

В отличие от классических подходов
HLBVH~\cite{HLBVH} и SBVH-presplitting~\cite{ganestam16_splitting},
работающих \emph{внутри} билдера BVH, предлагаемый метод оперирует
исключительно \emph{до} вызова вендорского BLAS-билдера, выполняется
на стороне runtime и явно учитывает специфический интерфейс DXR. Это позволяет получать ускорение, не имея
исходного кода драйвера и не модифицируя ассет-пайплайн движка.

\subsection{Практическая значимость}

Метод применим как опционально включаемый этап подготовки
динамических ассетов при инициализации BVH-контекста и не
требует переэкспорта ассетов. Он встраивается в существующий
рантайм-код движка без изменения формата хранения мешей, что
критично для проектов индустриального масштаба, использующих
большие массивы готовых ассетов.
